\subsection{稳定子群的拓扑直和}

由于带算子交换群的研究等同于模的研究，我们有时会容许自己对任意的带算子交换群使用专属于模的术语；这样，我们可以用线性映射一词来代替带算子交换群的同态，并可用投影算子一词来指带算子交换群的幂等自同态。

若带算子的交换拓扑群 \(\mathbf{E}\)（写作加法）是稳定子群的一个有限族 \((\mathbf{M}_i)_{1 \leqslant i \leqslant n}\) 的直和，则典范双射 \((x_i) \to \sum_{i=1}^{n} x_i\)（积群 \(\prod_{i=1}^{n} \mathbf{M}_i\) 到 \(\mathbf{E}\) 上的双射）是连续的，但未必是同胚。

定义 1. 设 E 是带算子的交换拓扑群，并设 \((\mathbf{M}_i)_{1 \leqslant i \leqslant n}\) 是 E 的稳定子群的一个有限族，使 E 是诸 \(\mathbf{M}_i\) 的直和。当下述条件成立时称 E 是诸 \(\mathbf{M}_i\) 的拓扑直和：典范映射 \((x_i) \to \sum_{i=1}^{n} x_i\)（积群 \(\prod_{i=1}^{n} \mathbf{M}_i\) 到 E 上的映射）是一个同胚（从而是带算子拓扑群的一个同构）。

命题 2. 设 E 是带算子的交换拓扑群，它是稳定子群 \(\mathbf{M}_i\)（ \(1 \leqslant i \leqslant n\)）的直和。设 \((p_i)_{1 \leqslant i \leqslant n}\) 是与分解 \(\mathbf{E} = \sum_{i=1}^{n} \mathbf{M}_i\) 相伴随的投影算子族。则 E 是诸 \(\mathbf{M}_i\) 的拓扑直和当且仅当诸 \(p_i\) 连续。

因为映射 \(x \to (p_i(x))\) 是映射 \((x_i) \to \sum_{i=1}^{n} x_i\) 的逆。

由于 \(\mathbf{I} = \sum_{i=1}^{n} p_i\)（其中 \(\mathbf{I}\) 表示 \(\mathbf{E}\) 的恒等映射），只需 \(n - \mathbf{I}\) 个投影算子 \(p_i\) 连续，便足以保证第 \(n\) 个也连续。

若 E 是两个稳定子群 M、N 的拓扑直和，则称 N 是 M 在 E 中的拓扑补；此事成立当且仅当 E/M 到 N 上的典范映射是带算子拓扑群的一个同构。

推论. 设 E 是带算子的交换拓扑群，并设 M 是 E 的稳定子群。那么下列条件等价：

\begin{enumerate}
\def\labelenumi{\alph{enumi})}
\item
  M 在 E 中有一个拓扑补。
\item
  存在连续投影算子 \(p\)（\(\mathbf{E}\) 到 \(\mathbf{E}\) 中的投影算子），使得 \(p(\mathbf{E}) = \mathbf{M}\)。
\item
  M 的恒等映射可以延拓为 E 到 M 上的一个连续线性映射。
\end{enumerate}

由命题 2 推出 a) 蕴含 b)，而 b) 蕴含 c) 是显然的。最后，若 p 是 E 到 M 上的一个连续线性映射，它延拓 M 的恒等映射，则 p 是一个连续投影算子，且投影算子 p 与 1-p 相伴于直和分解 \(E = M + N\)，其中 \(N = p^{-1}(\mathbf{o})\)。

注. 1) 为避免混淆，我们有时把 E 的作为 M 的补（就带算子群结构而言，不涉及拓扑）的稳定子群称为 M 的代数补。

\begin{enumerate}
\def\labelenumi{\arabic{enumi})}
\setcounter{enumi}{1}
\tightlist
\item
  若带算子的豪斯多夫交换拓扑群是稳定子群的一个族 \((\mathbf{M}_i)_{1\leqslant i\leqslant n}\) 的拓扑直和，则每个子群 \(\mathbf{M}_i\) 在 E 中闭，因为 \(\mathbf{M}_i\) 是由所有这样的 \(x\in E\) 组成的集合：它们满足 \(p_i(x) = x\)（第一章，§ 8，第 1 节，命题 2）。
\end{enumerate}

命题 3. 设 E、F 是两个带算子的交换拓扑群，并设 u 是 E 到 F 中的一个连续线性映射。为使存在 F 到 E 中的一个连续线性映射 v，使得 \(u \circ v\) 是 F 的恒等映射（这时称 u 为右可逆，称 v 为 u 的一个右逆），必须且只需 u 是 E 到 F 上的一个严格态射（§ 2，第 8 节），且 \(u^{-1}(\mathbf{o})\) 在 E 中有一个拓扑补。

这些条件是必要的。因为这时有 \(u(v(\mathbf{F})) = \mathbf{F}\)，并且 \(a\) fortiori（更有）\(u(\mathbf{E}) = \mathbf{F}\)；此外，若 \(p = v \circ u\)，则 \(p\) 是 \(\mathbf{E}\) 到其自身中的一个连续线性映射，满足 \(p^2 = p\)；因此（命题 2 的推论）\(p(\mathbf{E}) = v(u(\mathbf{E})) = v(\mathbf{F})\) 有一个拓扑补 \(p^{-1}(\mathbf{o})\)（在 \(\mathbf{E}\) 中的）；但按假设 \(u(p(x)) = u(x)\)，我们有 \(u^{-1}(\mathbf{o}) = p^{-1}(\mathbf{o})\)。最后，\(\mathrm{E}/p^{-1}(\mathbf{o})\) 到 \(\mathbf{F}\) 上的双射（与 \(u\) 相伴的双射），是如下两个映射的复合：\(\mathrm{E}/u^{-1}(\mathbf{o})\) 到 \(v(\mathbf{F})\) 上的双射（与 \(p\) 相伴的双射），以及 \(u\) 在 \(v(\mathbf{F})\) 上的限制；由于 \(v\) 连续，这两个映射都是同构，因此 \(u\) 是 \(\mathbf{E}\) 到 \(\mathbf{F}\) 上的一个严格态射。

这些条件是充分的。因为若 \(\varphi\) 是 E 到 \(\mathrm{E}/u^{-1}(\mathbf{o})\) 上的典范同态，则说 \(u^{-1}(\mathbf{o})\) 在 E 中有一个拓扑补 M，就等于说 \(\varphi\) 在 M 上的限制是 M 到 \(\mathrm{E}/u^{-1}(\mathbf{o})\) 上的一个同构。另一方面，由于 \(u=w\circ\varphi\)，其中 w 是 \(\mathrm{E}/u^{-1}(\mathbf{o})\) 到 F 上的一个同构，可见 u 在 M 上的限制是 M 到 F 上的一个同构，因而逆同构 v 使得 \(u\circ v\) 是 F 到其自身的恒等映射。

命题 4. 设 E、F 是两个带算子的交换拓扑群，并设 u 是 E 到 F 中的一个连续线性映射。为使存在 F 到 E 中的一个连续线性映射 v，使得 \(v \circ u\) 是 E 到其自身的恒等映射（这时称 u 为左可逆，称 v 为 u 的一个左逆），必须且只需 u 是 E 到 u(E) 上的一个（拓扑）同构，且 u(E) 在 F 中有一个拓扑补。

这些条件是充分的，因为若它们满足，就得到一个左逆 \(v\)（\(u\) 的左逆），其做法是把 \(u(\mathbf{E})\) 到 \(\mathbf{E}\) 上的同构（即 \(u\) 的逆）与 \(\mathbf{F}\) 到 \(u(\mathbf{E})\) 上的一个连续投影算子复合。

这些条件是必要的。因为关系 \(v(u(x)) = x\) 表明 \(u^{-1}(\mathbf{o}) = \{0\}\)；于是 \(u\) 是 \(E\) 到 \(u(E)\) 上的一个双射，且由于 \(v\) 在 \(u(E)\) 上的限制连续，由此得出 \(u\) 是 \(E\) 到 \(u(E)\) 上的一个同构。另一方面，若令 \(q = u \circ v\)，则 \(q\) 是 \(F\) 到 \(u(E)\) 上的一个连续线性映射，满足 \(q^2 = q\)，这就证明（命题 2 的推论）\(u(E)\) 在 \(F\) 中有一个拓扑补。

